\documentclass[tikz,border=4pt]{standalone}
\usepackage{ctex}
\usepackage{amsmath,amssymb}
\usetikzlibrary{arrows.meta,calc,decorations.markings}
\begin{document}
\begin{tikzpicture}[>=Stealth, line join=round,
  decoration={markings, mark=at position 0.55 with {\arrow{Stealth}}}
]
\draw[->] (-0.5,0) -- (5.5,0) node[right] {$x$};
\draw[->] (0,-0.5) -- (0,4.5) node[above] {$y$};

% 区域 D（椭圆形）
\fill[blue!15]
  plot[domain=0:360, samples=60]
    ({2.5+2*cos(\x)},{2+1.5*sin(\x)}) -- cycle;
\node[blue!70!black, font=\small] at (2.5,2) {$D$};

% 边界曲线 C（逆时针）
\draw[red, very thick, postaction={decorate}]
  plot[domain=0:360, samples=60]
    ({2.5+2*cos(\x)},{2+1.5*sin(\x)}) -- cycle;
\node[red, font=\small] at (4.9,2) {$C$（逆时针 $+$）};

% 箭头方向
\node[red, font=\tiny] at (2.5,3.7) {$\circlearrowleft$};

% 内部点和法向量/旋度
\foreach \xp/\yp in {1.5/1.5, 2.5/2, 3.5/1.5, 2/2.8, 3/2.8} {
  \draw[->, green!60!black, thick]
    (\xp,\yp) -- ++(-0.15,0.4)
    node[above, font=\tiny] {};
}
\node[green!50!black, font=\small] at (1.0,3.0) {$Q_x-P_y$（旋度密度）};

% Green 公式标注
\node[font=\small, draw=gray!50, fill=gray!10, rounded corners=3pt, align=center]
  at (2.5,-0.3)
  {\textbf{Green 定理}：$\oint_C P\,dx+Q\,dy = \iint_D\!\!\left(\frac{\partial Q}{\partial x}-\frac{\partial P}{\partial y}\right)\!dA$};

% 图例
\draw[red, very thick] (0.5,4.2) -- (1.1,4.2) node[right, font=\small, black] {边界 $C$（逆时针）};
\fill[blue!15] (0.5,3.8) rectangle (1.1,4.0);
\draw[blue!60] (0.5,3.8) rectangle (1.1,4.0);
\node[right, font=\small] at (1.1,3.9) {区域 $D$};
\end{tikzpicture}
\end{document}
